\documentclass[11pt]{article}
\usepackage[margin=1in]{geometry}
\usepackage{amsmath,amssymb}
\usepackage{xcolor}
\usepackage[hidelinks]{hyperref}
\usepackage{enumitem}
\setlist{leftmargin=*,itemsep=0.5em}
\newcommand{\manuscript}{\textit{Construction of Kochen--Specker Sets from Mutually Unbiased Bases}}

\begin{document}

\begin{center}
{\Large Response to the Editor}\\[0.5em]
\manuscript\\[0.25em]
Mirko Navara and Karl Svozil\\[0.25em]
4 August 2026
\end{center}

Dear Editor,

Thank you for the opportunity to revise our manuscript.  We appreciate the first referee's favorable assessment and recommendation to publish the paper without revision.  We also carefully considered the second referee's request for a substantial presentational revision.  We found those requests reasonable and have implemented them in full.  All additions and changed passages are shown in blue in the revised manuscript.

The revision has two parts.  First, we addressed every organizational concern raised by the second referee:

\begin{itemize}
\item the Introduction is shorter and more focused;
\item rays, contexts, intertwining rays, two-valued states, KS configurations, and maximal unbiasedness are defined when first used;
\item the main claims are isolated as three propositions with proofs or exact verification procedures;
\item the qutrit construction now begins with a four-stage guide;
\item the $D=4$ and $D=5$ forcing constraints are separated from coordinate completions and computational inventories; and
\item repetitive prose and speculative wording have been removed or moderated.
\end{itemize}

Second, we performed an independent exact audit of the mathematical inventories and the higher-dimensional comparison.  This led to several important corrections:

\begin{enumerate}
\item The claimed 40--4--4 provenance asymmetry was a labeling artifact.  The invariant decomposition is symmetric: 21 rays and 10 contexts are common to all three qutrit lineages; each lineage adds 48 exclusive rays and 40 lineage-specific contexts.  Thus every lineage has 69 rays and 50 contexts, and their union has 165 rays and 130 contexts.

\item The concrete 20-ray $D=4$ set supports exactly 17 orthogonal contexts, rather than the previously reported 13.  Its complete context hypergraph has 16 separating two-valued states.  The revised paper also distinguishes the selected 10- and 13-context subhypergraphs, both of which have the same 36 states, from the complete 17-context hypergraph.

\item The comparison with the Cabello 18-ray set has been corrected.  The two sets share 15 rays projectively.  Five gadget-only rays and three Cabello-only rays are mutually reconstructible by rank-three orthogonal complements.  We corrected the noncollinear $v_{1423}$/$v_{47}$ pairing and three invalid reconstruction triples.

\item We removed a four-dimensional MUB appendix whose displayed bases did not satisfy the claimed mutual-unbiasedness condition.  The higher-dimensional forcing proofs need only the verified equal-modulus criterion, which is now stated directly; the appendix is restricted to the qutrit MUBs actually used in the construction.

\item We corrected the logical explanation of context completion.  Adding contexts can only preserve or reduce the set of two-valued states; it cannot create ``loopholes.''  The revised manuscript compares distinct ray/context selections inside the common Peres--Mermin parent without violating this monotonicity.
\end{enumerate}

The revised source compiles cleanly, with resolved citations and cross-references, and the resulting ten-page PDF has been checked page by page.  We believe the manuscript is now substantially clearer and technically stronger, and we respectfully submit it for reconsideration.

Sincerely,

Mirko Navara and Karl Svozil

\end{document}
